\section{Ordinary symmetric-square congruences in every weight of a progression}
\label{sec:kt-ordinary-symmetric-square}

Put $E=\mathbf Q(\sqrt{-3})$ and $\epsilon=\chi_{-3}$.
The cyclotomic character $\varepsilon_\ell$ below is normalized
by $\varepsilon_\ell(\operatorname{Frob}^{\rm arith}_q)=q$.
This section strengthens KT-U3-SYM2-EISENSTEIN by verifying
the ordinary local flags and characteristic-zero irreducibility.
Its last source-predicate implication is stated separately so
that it can be checked against the literal definitions.

\begin{theorem}[KT-ORDINARY-SYM2-FAMILY]
Let $e_0\geq2$ be even and let $\ell>3$ be prime. Suppose
\[
 v_\ell\!\left(\frac{B_{e_0+1,\epsilon}}{e_0+1}\right)>0,
 \qquad \ell\nmid1-2^{e_0}.                         \tag{OS.1}
\]
For each $t\geq0$ put $e=e_0+(\ell-1)t$. There is a
non-CM primitive normalized eigenform
\[
 f_t\in S_{e+1}(\Gamma_0(3),\epsilon)
\]
and a coefficient prime $\lambda\mid\ell$ such that
\[
 a_q(f_t)\equiv1+\epsilon(q)q^e\pmod\lambda
                       \quad\text{for every prime }q\ne3.
                                                               \tag{OS.2}
\]
In particular $a_\ell(f_t)$ is a $\lambda$-adic unit.
Let $r_f$ be its arithmetic-normalized two-dimensional
$\ell$-adic representation and set
\[
 \rho_t=\left(\operatorname{Sym}^2r_f\otimes
                         (\det r_f)^{-1}\right)|_{G_E}.
                                                               \tag{OS.3}
\]
Then $\rho_t$ has all of the following properties.
\begin{enumerate}
\item It is continuous and absolutely irreducible over
      $\overline{\mathbf Q}_\ell$, with determinant one.
\item It is unramified at every finite place of $E$ away from
      $\ell$. At every place above $\ell$ it has an actual
      upper triangular flag, in the displayed order, with
      inertia characters
      $\varepsilon_\ell^e,1,\varepsilon_\ell^{-e}$.
\item For the three diagonal characters $\kappa_i$ on any
      such decomposition group and all $i,j$ and $m\geq1$,
      $\kappa_i^m\ne(\kappa_j\varepsilon_\ell)^m$.
      The same assertion holds after twisting the whole
      representation by $\varepsilon_\ell^{-1}$.
\item It is the weight-zero system of the everywhere
      finite-unramified, regular algebraic, conjugate
      self-dual cuspidal representation
      $\Pi_f=\operatorname{BC}_{E/\mathbf Q}
                    (\operatorname{Ad}(\pi_f))$.
      The normalized system used in \texttt{AttC} is
      $\rho_t\varepsilon_\ell^{-1}$.
\item For a stable lattice its residual semisimplification is
\[
       \overline\rho_t^{\,\rm ss}
           \simeq\overline\varepsilon_\ell^e\oplus1\oplus
                           \overline\varepsilon_\ell^{-e}
                      \quad\text{on all of }G_E.              \tag{OS.4}
\]
\end{enumerate}
Consequently $\rho_t$ satisfies the literal predicate
$\mathrm{Level}(\ell,2,(e,0,-e))$, including its ordinary
flag and nonadjacency alternative. The representation is
also a stable compact $U_3$ eigenline in the constructed
symmetric-square subspace of KT-U3-SYM2-INJECT.
\end{theorem}

\begin{proof}
\paragraph{The infinite progression and ordinarity.}
Write $w=e+1$ and
$\beta_w=B_{w,\epsilon}/w$. Since $w$ is odd and $\ell-1$
is even, $\ell-1\nmid w$. Sun, \emph{Congruences concerning
Bernoulli numbers and Bernoulli polynomials}, Lemma~8.1(a),
printed pp.~219--220, applied with $n=1$, gives
\[
 (1-\epsilon(\ell)\ell^{w-1})\beta_w
 \equiv
 (1-\epsilon(\ell)\ell^{w+\ell-2})\beta_{w+\ell-1}
                         \pmod\ell.                         \tag{OS.5}
\]
The conductor is three, prime to $\ell$. Each Euler factor
is an $\ell$-adic unit since $w\geq3$. Starting at (OS.1)
and applying (OS.5) successively proves $v_\ell(\beta_w)>0$
for every weight in the progression. The lemma concerns
the normalized Bernoulli numbers themselves and remains
valid when $w$ is divisible by $\ell$.

Fretwell--Roberts, \emph{Newform Eisenstein congruences of
local origin}, Theorem~3.6, printed pp.~517--518, now applies
with $N=3$, $M=1$, $k=w$, $\psi=1$, $\phi=\epsilon$.
Its special value is $L(1-w,\epsilon)=-\beta_w$;
its coefficient-prime exclusion is just $\ell\nmid6$.
There is no hypothesis $\ell>k+1$ in this sufficient theorem.
It gives one eigenform and one coefficient prime satisfying
(OS.2) simultaneously for every prime $q\nmid NM$.
In particular $q=\ell$ is included, so $a_\ell\equiv1$.
Primitivity of the nebentypus makes the eigenform new, as
explained immediately after that theorem's proof.

All CM forms in this level-three space are CM by $E$, by
KT-GL2-LOWER-SCOPE. At the inert prime two they have $a_2=0$.
But Fermat's theorem and (OS.1) give
$a_2\equiv1-2^e\equiv1-2^{e_0}\ne0\pmod\lambda$.
Thus every chosen $f_t$ is non-CM.

\paragraph{The ordinary flag, with its orientation.}
The ordinary local theorem of Mazur--Wiles and Wiles says
that $a_\ell$ a unit gives an upper triangular realization
\[
        r_f|_{G_{\mathbf Q_\ell}}\sim
           \begin{pmatrix}\delta&*\\0&\eta\end{pmatrix},
       \qquad \eta\text{ unramified},\qquad
       \delta\eta=\epsilon\varepsilon_\ell^e.                \tag{OS.6}
\]
An exact published statement of the general theorem, with
this orientation, is Ghate--Vatsal, \emph{On the local
behaviour of ordinary $\Lambda$-adic representations},
printed p.~2144, equation~(1.1) and its preceding paragraph.
No residual irreducibility assumption is imposed on this
ordinary assertion. Their later splitting theorems impose
additional assumptions; none is used here.

If $v_1,v_2$ is the basis in (OS.6), the ordered basis
$v_1^2,v_1v_2,v_2^2$ gives an upper triangular symmetric
square with diagonal $\delta^2,\delta\eta,\eta^2$.
Dividing by $\delta\eta$ yields the actual flag of (OS.3)
with diagonal
\[
                    \delta/\eta,\quad1,\quad\eta/\delta.
                                                               \tag{OS.7}
\]
Because $\ell\ne3$, the character $\epsilon$ is unramified
at $\ell$; in fact its restriction to $G_E$ is trivial.
On inertia (OS.7) is exactly
$(\varepsilon_\ell^e,1,\varepsilon_\ell^{-e})$.
Restriction to each completion of $E$ preserves the flag.
Conjugation transports it to every choice of a place in the
algebraic closure; no splitting of the flag is asserted.

Let $h=(e,0,-e)$. If
$\kappa_i^m=(\kappa_j\varepsilon_\ell)^m$ for $m\geq1$,
restriction to inertia gives
$\varepsilon_\ell^{m(h_i-h_j-1)}=1$.
The cyclotomic image of inertia in a finite extension of
$\mathbf Q_\ell$ is infinite. Hence $h_i-h_j=1$ would
follow. This is impossible since
$h_i-h_j\in\{0,\pm e,\pm2e\}$ and $e\geq2$.
Twisting all characters by $\varepsilon_\ell^{-1}$ leaves
their differences unchanged. This proves every asserted
positive-power nonlink, not just the case $m=1$.

\paragraph{Characteristic-zero irreducibility at every prime.}
Ribet, \emph{On $\ell$-adic representations attached to modular
forms II}, printed p.~191, immediately before Theorem~3.1,
recalls Momose's theorem that, for each $\ell$, the image
on the fixed inner-twist-kernel subgroup is open in the
reduced-norm group of the specified quaternion algebra.
The original reference is F.~Momose, \emph{On $\ell$-adic
representations attached to modular forms}, J.~Fac.~Sci.\
Univ.~Tokyo Sect.~IA Math.~28 (1981), 89--109.
This is the all-prime input used here; Ribet's Theorem~3.1
itself is an almost-all-primes statement and is not used
as an all-prime theorem.
After extension of coefficients to $\overline{\mathbf Q}_\ell$,
the projection to the chosen coefficient embedding contains
$\mathfrak{sl}_2$. This follows as well by taking the
commutator of the central simple algebra in that theorem.
Restriction to the finite-index subgroup $G_E$ does not
change the Lie algebra of the image. The representation
$\operatorname{Sym}^2\otimes\det^{-1}$ of
$\mathfrak{sl}_2$ is its irreducible three-dimensional
adjoint representation in characteristic zero. Therefore
any invariant subspace of (OS.3) is zero or the whole space.
This proves absolute irreducibility even when the residual
representation is reducible. Ribet's separate almost-all-
primes integral-image theorem is not substituted for this
all-prime Lie algebra theorem.

\paragraph{The global and residual assertions.}
KT-U3-SYM2-INJECT applies to every even $e\geq2$ and every
non-CM $f_t$. It proves the cuspidal base change, its regular
algebraic conjugate self-duality, all finite unramified local
parameters, and the stable compact member. Its local check
at three uses inertia $1\oplus\epsilon$ and zero monodromy
for the classical form; restriction to $E_3$ kills
$\epsilon$. Classical local compatibility therefore makes
(OS.3) unramified at three as well as at every other place
away from $\ell$.

The weight-zero adjoint is matched to $\Pi_f$ by its good
Frobenius polynomials; its three Satake roots are invariant
as a multiset under inversion. The definition \texttt{AttC} uses
geometric Frobenius and the factor $q^{(3-1)/2}=q$ in
\texttt{algSat}; twisting the weight-zero system by
$\varepsilon_\ell^{-1}$ supplies exactly this factor.
Thus this is the automorphic normalization required by
\texttt{Adm3Alt}.

Continuity and a finite coefficient field give a stable
lattice. The determinant of $r_f$ is
$\epsilon\varepsilon_\ell^e$. The simultaneous congruence
(OS.2), Chebotarev, and Brauer--Nesbitt give
\[
       \overline r_f^{\,\rm ss}
              \simeq1\oplus\epsilon\overline\varepsilon_\ell^e.
\]
Restriction to $G_E$ kills $\epsilon$. Taking the symmetric
square and dividing by the determinant gives the three
composition factors in (OS.4); equivalently, use the flag
of the endomorphisms and remove the scalar line. Since
$\ell>3$, scalar and trace-zero endomorphisms split.

Finally, at level parameter $2$ choose all the global finite
characters in \texttt{LevelAway} and \texttt{LevelAt} to be
trivial. They satisfy \texttt{KillsMu 2}. Unramifiedness
gives the required diagonal inertia away from $\ell$;
(OS.7) gives \texttt{IsFlagOn} above $\ell$ and the required
inertia powers in the right order. The case $m=1$ of the
nonlink argument supplies the first disjunct of
\texttt{LevelAt}. No potentially crystalline alternative or
Fontaine--Laffaille range is needed.
\end{proof}

\begin{corollary}[KT-ORDINARY-SYM2-VANISHING-PROGRESSIONS]
Under the stipulated archive-theorem convention used in
KT-K52-K56-K60-VANISH, for every row
\[
\begin{array}{c|r}
e_0&M_{e_0}\\\hline
26&697446897311\\
28&323026426255399\\
30&236017358500033404991
\end{array}                                                   \tag{OS.8}
\]
and every prime $\ell\mid M_{e_0}$, all the conclusions of
KT-ORDINARY-SYM2-FAMILY hold for
$e=e_0+(\ell-1)t$, $t\geq0$. In addition,
\[
                  K_{2e}(\mathbf Z)\{\ell\}=0
                           \qquad(t\geq0).                   \tag{OS.9}
\]
\end{corollary}
\begin{proof}
The exact values in KT-U3-SYM2-EISENSTEIN, or the independent
rational-arithmetic certificate
\texttt{lower\_bernoulli\_certificate.json}, give
$M_{e_0}\mid\operatorname{num}(B_{e_0+1,\epsilon}/(e_0+1))$
and
$\gcd(M_{e_0},6\operatorname{den}(B_{e_0+1,\epsilon}/(e_0+1))
 (1-2^{e_0}))=1$.
Thus each prime divisor satisfies (OS.1) and is greater than
three. No primality assumption about the displayed moduli
is required.

KT-K52-K56-K60-VANISH gives $K_{2e_0}(\mathbf Z)=0$.
The carried comparison (K2), with its all-degree scope
K17-K2-PERIODIC, identifies the vanishing criterion with the
same cyclotomic class-group character indexed by $-e$.
Since $e\equiv e_0\pmod{\ell-1}$, that character and hence
the criterion are unchanged. This proves (OS.9). The
argument uses periodicity of the vanishing criterion, not
an isomorphism between $K$-groups in different degrees.
\end{proof}

\begin{corollary}[KT-ORDINARY-SYM2-ADM3]
Under the stipulated archive-theorem convention, the
representations in KT-ORDINARY-SYM2-FAMILY satisfy the actual
predicate
\[
                 \mathrm{Adm3U\_s}(\iota,e,2,\rho_t).          \tag{OS.10}
\]
Here $\iota:\overline{\mathbf Q}_\ell\simeq\mathbf C$ is
chosen compatibly with the classical eigenform and its
attached automorphic representation. This includes the
stronger \texttt{WitA} and the added \texttt{Th.Unl}
conditions, not just ordinary automorphy and level.
\end{corollary}
\begin{proof}
At the pinned source revision
\texttt{7cbf7a0c32fca1badf05931879a37c3f8a2b03b4},
\texttt{Adm3Alt} requires \texttt{IsIrr},
\texttt{IsAutCsdA} for $\rho\varepsilon_\ell^{-1}$, and
the precise \texttt{Level} proved above.
\texttt{Adm3U\_s} adds \texttt{Th.Unl}.
Write $r=\rho_t\varepsilon_\ell^{-1}$ and let
$\pi_0=\Pi_f$. We have \texttt{AutCsd0}$(\iota,r)$ by the
regular algebraic conjugate self-dual cuspidal lift and
the normalized attachment already proved.

\paragraph{The first clause of \texttt{WitA}.}
The diagonal inertia exponents of $r$ are
$(e-1,-1,-e-1)$. The proved positive-power nonlink gives
\texttt{PotUnl}$(r,v)$ at every $v\mid\ell$, for every
choice of the valuation above $v$. It therefore proves,
in particular, the implication
\texttt{Unl1}$(r,v)\Rightarrow$\texttt{PotUnl}$(r,v)$
required by the first clause of \texttt{WitA}. The closed
source theorem \texttt{potOfFlag} gives exactly the same
deduction: take its exponent tuple to be
$(e-1,-1,-e-1)$ and its auxiliary power to be one.

\paragraph{The universally quantified window clause.}
Every finite local component of $\pi_0$ is spherical.
The closed theorem \texttt{sphUnl} consequently supplies
toral support and unlinked characters at every such place.
The closed theorem \texttt{winOf}, applied once to $\pi_0$,
gives a single real number $u$ such that
\[
                  \mathrm{WinAt}(\pi_0,u,v)
                        \quad\text{for every finite }v.     \tag{OS.11}
\]
The same $u$ works at every place: the existential real
parameter in \texttt{WinOf} precedes its place quantifier.
This uses the exact strict width-one window of the source.
It does not require an additional Ramanujan assertion.

Now let $\pi$ be any regular algebraic cusp representation
with \texttt{AttC}$(\iota,r,\pi)$. Its attachment and that of
$\pi_0$ give the same characteristic polynomial at all but
finitely many places. Both use the same nonzero $q$-factor
in \texttt{algSat}; uniqueness of the roots with
multiplicity gives the same Satake multiset there.
The closed theorem \texttt{suppSmo}, whose proof applies
\texttt{strongMultOneC}, transports (OS.11) to $\pi$.
Thus the same window holds at every place for this $\pi$,
and in particular at every place selected by the disjunction
in the second clause of \texttt{WitA}. This proves its
universal quantifier over attached cusp representations.

Since \texttt{IsAutCsdA} is precisely \texttt{AutCsd0}
together with the implication from continuity to
\texttt{WitA}, the preceding paragraphs prove that
predicate for $r$.

\paragraph{The extra unlinked flag and conclusion.}
The definition \texttt{FlagAt} asks for the flag, the
inertia powers and global finite characters killed by
$\mu_2$. Choose those characters to be one and use (OS.7).
The case $m=1$ of the proved nonlink is the remaining
inequality in \texttt{Th.Unl}. Thus \texttt{Th.Unl}
holds for $\rho_t$ and $(e,0,-e)$.
Combining this with absolute irreducibility,
\texttt{IsAutCsdA} for $r$, and the proved \texttt{Level}
constructs \texttt{Adm3Alt} and then
\texttt{Adm3U\_s}, exactly as defined.
\end{proof}

\begin{corollary}[KT-ORDINARY-SYM2-GD]
For every row of (OS.8), every prime $\ell\mid M_{e_0}$
and every $t\geq0$, the parameters $E$, $N=2$, trivial
Dirichlet character $\psi$, and $e=e_0+(\ell-1)t$ satisfy
the actual predicate \texttt{GD} in \texttt{OneCharU\_s}.
\end{corollary}
\begin{proof}
The already checked rational specialization in
\texttt{SmallClassNumbers/ClassNumberThree.lean},
lines 64--86, proves \texttt{GD} for every positive even
$e$ and prime $\ell\geq5$ with $(\ell-1)\nmid2e$.
It has no additional inequality comparing $e$ with $\ell$.
For the three rows of (OS.8), exact integer arithmetic gives
\[
                   \gcd\bigl(M_{e_0},(2e_0+1)!\bigr)=1.
\]
Thus each prime divisor satisfies $\ell-1>2e_0>0$.
Since $2e\equiv2e_0\pmod{\ell-1}$, its nondivisibility
condition persists for every $t$. The exact gcd outputs
are saved in \texttt{ordinary\_family\_gd\_certificate.json}.
\end{proof}

\begin{corollary}[KT-ORDINARY-SYM2-EXPLICIT-LARGE]
The integer $\ell=697446897311$ is prime and exceeds
$2^{31}$. For every $t\geq0$, put
\[
                  e=26+697446897310\,t.
\]
Then the actual \texttt{GD} parameters of
KT-ORDINARY-SYM2-GD admit a representation satisfying
\texttt{Adm3U\_s}, the global residual triple (OS.4),
and $K_{2e}(\mathbf Z)\{\ell\}=0$.
Thus full-field unit exclusion fails at a fixed prime
above $2^{31}$ in unbounded weights within the literal
source hypotheses of \texttt{OneCharU\_s}.
\end{corollary}
\begin{proof}
All assertions except primality follow from the preceding
three corollaries. Here is an exact primality certificate
independent of trial division of the target $\ell$.
Write $n=697446897311$. Then
\[
              n-1=2\cdot5\cdot31\cdot1019\cdot2207879.
\]
Each displayed factor is prime: exhaustive division up to
its integer square root suffices, and the largest such
bound is $1485$. For the following witnesses $a_q$, exact
modular exponentiation gives $a_q^{n-1}\equiv1\pmod n$
and the displayed $b_q=a_q^{(n-1)/q}\bmod n$:
\[
\begin{array}{r|r|r}
q&a_q&b_q\\\hline
2&13&697446897310\\
5&2&682051080219\\
31&2&34767638009\\
1019&2&273082850090\\
2207879&2&201208097879
\end{array}
\]
In each row $\gcd(b_q-1,n)=1$.
For any prime $p\mid n$, these conditions imply that the
order of $a_q$ modulo $p$ contains the full $q$-power of
$n-1$. Since this holds for every prime $q\mid n-1$,
we obtain $n-1\mid p-1$. Thus $p\geq n$, forcing $p=n$.
This proves primality. The inequality $n>2^{31}$ is
immediate (the difference is $695299413663$).
The complete standard-library verification is
\texttt{ordinary\_large\_prime\_certificate.py}, with
exact output saved beside it.
\end{proof}

\paragraph{Exact source locators for (OS.10).}
All paths here are relative to
\texttt{src/Development/AutomorphicArithmetic/}.
The predicates \texttt{IsFlagOn}, \texttt{LevelAway},
\texttt{LevelAt}, and \texttt{Level} are in
\texttt{AutomorphicVarieties/AutomorphicForms/%
BlochKatoSelmerVanishes.lean}, lines 37--39,
89--92, 102--107, and 111--114;
\texttt{FlagAt} is in the same directory's
\texttt{LocIsotypic.lean}, lines 31--34.
\texttt{AutCsd0} is in
\texttt{Foundations/GaloisRepresentations/AdmissibleResidues/%
IsAutomorphicOverImagQuad.lean}, lines 35--36.
Inside \texttt{AutomorphicVarieties/ModularityArgument/},
\texttt{AutomorphyLifting/WeakOfAtt.lean}, lines 19--30,
defines \texttt{Unl1}, \texttt{PotUnl}, and \texttt{WitA};
\texttt{AutomorphyLifting/IsAutCsdA.lean}, lines 11--12,
defines \texttt{IsAutCsdA}.
The definitions \texttt{WinAt}, \texttt{SuppSmo},
\texttt{WinOf}, and \texttt{SphUnl} are respectively at
lines 38--41, 142--145, 158--164 of
\texttt{LeviSmoothLocalization/WinUpDown.lean}.
In its \texttt{PrincipalSeriesMatching/} subdirectory,
\texttt{PotOfFlag.lean}, lines 88--99,
\texttt{SphUnl.lean}, lines 25--45,
\texttt{NormAbsPowUnif.lean}, lines 124--168, and
\texttt{SuppSmo.lean}, lines 6--12, contain the closed
theorems just used, with their complete proofs.
Finally \texttt{AutomorphicVarieties/TheoremsMuLinked/%
AutP.lean}, lines 10--13,
defines \texttt{Th.Unl}; the
\texttt{ModularityArgument/ThresholdArguments/%
DenseThresholdCases/} files \texttt{AfStmt.lean},
lines 29--32, and \texttt{NumberU3US.lean}, lines 27--29,
define \texttt{Adm3Alt} and \texttt{Adm3U\_s}.

\paragraph{Consequence and scope.}
For every representation in (OS.10), (OS.4) makes
\[
 \operatorname{tr}\rho_t(\sigma)
  -\bigl(\varepsilon_\ell(\sigma)^e+1+
                         \varepsilon_\ell(\sigma)^{-e}\bigr)
\]
belong to the coefficient maximal ideal for every
$\sigma\in G_E$. Thus varying a single group element,
or any number of good-prime Hecke operators, cannot make
this trace difference a unit for all members of the
actual admissible rank-three comparison class.
The families (OS.8)--(OS.9) show that this obstruction
coexists with the vanishing of the indicated rational
$K$-group component.
Membership in \texttt{Adm3U\_s} is not by itself a converse
to the source's map from deformation-ring points.
The next section supplies the additional argument:
KT-ADJOINT-ACTUAL-HOST-POINT realizes these representations
as actual Type-I and TwoPt points, so every literal
\texttt{CoversPt} cover contains their isomorphism class.
This stronger conclusion still does not assert that every
possible existential witness set in \texttt{OneCharU\_s}
must contain them.

\paragraph{Primary inputs and check scope.}
Sun's Lemma~8.1(a), including its displayed proof, and
Fretwell--Roberts Theorem~3.6, including its complete proof,
were already read and independently checked in the preceding
family audit. The present ordinary input is checked at the
exact published statement of Ghate--Vatsal, p.~2144,
equation~(1.1); the original Wiles 1988 proof is an imported
theorem and is not claimed to have been reread. The exact
Momose open-image statement recalled by Ribet on p.~191
is the imported all-prime theorem; no reread of Momose's
original proof is claimed. An earlier draft incorrectly
called this Ribet Theorem~0.1 and was corrected after the
independent visual source check. The global lift is the independently
checked current-run KT-U3-SYM2-INJECT. Source PDFs and
hashes for the two new local/irreducibility inputs are
saved under \texttt{evidence/main\_gate/literature/}.
